%%%%%%%%%%%%%%%%%%%%%%%%%%%%%%%%%%%%%%%%%%%%%%%%%%%%%%%%%%%%%%%%%%%%%%
%%  Copyright by Wenliang Du and SEED Lab contributors.             %%
%%  This work is licensed under the Creative Commons                %%
%%  Attribution-NonCommercial-ShareAlike 4.0 International License. %%
%%  To view a copy of this license, visit                           %%
%%  http://creativecommons.org/licenses/by-nc-sa/4.0/.              %%
%%%%%%%%%%%%%%%%%%%%%%%%%%%%%%%%%%%%%%%%%%%%%%%%%%%%%%%%%%%%%%%%%%%%%%

\newcommand{\commonfolder}{../../../common-files}

\input{\commonfolder/header}
\input{\commonfolder/copyright}
\hypersetup{pdfborder = {0 0 0}}

\lhead{\bfseries SEED Labs -- BGP RPKI and ROV Lab}

\begin{document}

\begin{center}
{\LARGE BGP RPKI and Route Origin Validation Lab}
\end{center}

\seedlabcopyright{2026}


% *******************************************
% SECTION
% *******************************************
\section{Overview}

Border Gateway Protocol (BGP) allows autonomous systems (ASes) to announce
which IP prefixes they can reach. This design gives the Internet its global
connectivity, but it also creates a serious security problem: BGP speakers
traditionally trust route announcements from their neighbors. If an AS announces
a prefix that it does not own, other networks may route traffic to the wrong
place. Such incidents can be caused by attacks, operator mistakes, or route
leaks.

The Resource Public Key Infrastructure (RPKI) helps address this problem. In
RPKI, a prefix holder can create a Route Origin Authorization (ROA), which says
which AS is authorized to originate a prefix and what maximum prefix length is
allowed. Routers can then perform Route Origin Validation (ROV) on received BGP
routes. A route can be classified as \textit{Valid}, \textit{Invalid}, or
\textit{NotFound}. Operators can reject Invalid routes to reduce the impact of
origin hijacking.

The objective of this lab is to help students understand how RPKI and ROV work
and what they do not protect against. Students will use the SEED Internet
Emulator to launch a prefix hijacking attack, create a small educational ROA
table, install a BIRD import filter that models ROV, and show that the
unauthorized route is rejected. Students will also test a more-specific prefix
attack and discuss the limitations of origin validation.

This lab covers the following topics:

\begin{itemize}[noitemsep]
\item BGP prefix announcements and origin ASes
\item Prefix hijacking and more-specific prefix attacks
\item RPKI, ROA, maxLength, and Route Origin Validation
\item BIRD route filters in the SEED Internet Emulator
\item Security limitations of ROV
\end{itemize}

\paragraph{Readings.}
Students should read the SEED BGP tutorial and the existing BGP Exploration and
Attack Lab before working on this lab. The following public standards are also
useful:

\begin{itemize}
\item RFC 6480, An Infrastructure to Support Secure Internet Routing:
      \url{https://www.rfc-editor.org/rfc/rfc6480}
\item RFC 6811, BGP Prefix Origin Validation:
      \url{https://www.rfc-editor.org/rfc/rfc6811}
\item RFC 8210, The RPKI to Router Protocol:
      \url{https://www.rfc-editor.org/rfc/rfc8210}
\item NIST NCCoE RPKI project page:
      \url{https://www.nccoe.nist.gov/projects/routing-security}
\end{itemize}

\paragraph{Lab environment.}
\seedenvironmentB
\nodependency


% *******************************************
% SECTION
% *******************************************
\section{Lab Environment}

This lab is conducted inside the SEED Internet Emulator. We use the same
emulator conventions as the BGP Exploration and Attack Lab:

\begin{itemize}[noitemsep]
\item AS numbers \texttt{150}--\texttt{199} are stub ASes.
\item The first network prefix for AS-\texttt{N} is \texttt{10.N.0.0/24}.
\item BGP routers run BIRD, and their configuration file is
      \path{/etc/bird/bird.conf}.
\item The emulator map can be accessed at \url{http://localhost:8080/map.html}
      after the emulator starts.
\end{itemize}

\paragraph{Starting the emulator.}
Download the BGP emulator lab setup file, unzip it, and go to the generated
\texttt{output} folder. Then build and start the containers:

\begin{lstlisting}
$ docker-compose build
$ docker-compose up

// Aliases in the SEED VM
$ dcbuild
$ dcup
$ dcdown
\end{lstlisting}

\paragraph{Useful commands.}
The following commands are used throughout the lab:

\begin{lstlisting}
// List containers
$ dockps

// Get a shell inside a container
$ docksh <container-id>

// Show BIRD protocols
# birdc show protocols

// Show all BGP routes
# birdc show route all

// Reload BIRD configuration after editing /etc/bird/bird.conf
# birdc configure
\end{lstlisting}

\paragraph{Helper files.}
This lab folder contains a \texttt{Labsetup} folder with an educational ROA
table and a helper script:

\begin{lstlisting}
Labsetup/
|-- roa-table.txt
|-- generate_bird_rov_filter.py
\end{lstlisting}

The helper script generates a BIRD filter function from the ROA table. This is
not a full RPKI validator and it does not use the RTR protocol. It is a small
model that lets students learn the policy logic behind ROV inside the emulator.


% *******************************************
% SECTION
% *******************************************
\section{How RPKI and Route Origin Validation Work}

Before working on the lab tasks, we first describe the theory behind RPKI and
Route Origin Validation. BGP routers exchange reachability information, but
traditional BGP does not include a built-in way to prove that an AS is allowed
to originate a prefix. If AS-161 announces \texttt{10.154.0.0/24}, a normal
BGP router can evaluate path attributes and policies, but it cannot know from
BGP alone whether AS-161 is authorized to originate that prefix.

\subsection{Resource Certificates}

RPKI is based on the same general idea as public-key infrastructure, but the
certificates bind Internet number resources to public keys. These resources
include IP address blocks and AS numbers. Regional Internet Registries (RIRs)
are trust anchors for the resources they allocate. A resource holder can receive
a certificate proving that it controls a particular prefix.

For this lab, we do not build the full certificate hierarchy. Instead, we model
the result of the hierarchy: a trusted database telling routers which ASes are
authorized to originate which prefixes.

\subsection{Route Origin Authorization}

A Route Origin Authorization (ROA) is a signed object created by a prefix
holder. It authorizes one AS to originate a prefix, optionally allowing more
specific prefixes up to a specified maximum length. A ROA contains three
important fields:

\begin{itemize}[noitemsep]
\item Prefix: the address block covered by the ROA.
\item Origin AS: the AS authorized to originate the prefix.
\item maxLength: the longest prefix length that is still authorized.
\end{itemize}

For example, the following educational ROA entry says that AS-154 can originate
\texttt{10.154.0.0/24}, but not a more-specific route such as
\texttt{10.154.0.0/25}.

\begin{lstlisting}
10.154.0.0/24,24,154,AS-154 customer prefix
\end{lstlisting}

If the maxLength were changed to \texttt{25}, then AS-154 would be authorized
to originate both \texttt{10.154.0.0/24} and covered \texttt{/25} prefixes.
Choosing maxLength is therefore a security decision. A loose maxLength can make
some accidental or malicious more-specific announcements look valid.

\subsection{Validation States}

When a router receives a BGP route, it can compare the route's prefix and origin
AS with the validated ROA database. The result is one of three states:

\begin{itemize}[noitemsep]
\item \textbf{Valid:} A covering ROA exists, the origin AS matches, and the
      prefix length is no longer than maxLength.
\item \textbf{Invalid:} A covering ROA exists, but the origin AS is wrong or
      the prefix length is longer than maxLength.
\item \textbf{NotFound:} No covering ROA exists for the prefix.
\end{itemize}

The origin AS is the last AS in the AS path. For example, if the path is
\texttt{3 161}, the origin AS is AS-161. Route Origin Validation checks this
origin AS. It does not validate every AS hop in the path.

\subsection{Router Policy}

ROV does not automatically force a router to reject a route. It gives the router
a validation result, and the operator decides the policy. A common security
policy is to reject Invalid routes while continuing to accept Valid and NotFound
routes. Operators often accept NotFound routes because not every legitimate
prefix on the Internet has a ROA.

In a real deployment, an RPKI validator downloads and validates signed RPKI
objects from repositories. It then sends validated prefix-origin information to
routers, commonly using the RPKI-to-Router (RTR) protocol. In this lab, we use a
small ROA table and a generated BIRD filter to model the same decision:

\begin{center}
BGP route + ROA table $\rightarrow$ validation state $\rightarrow$ routing policy
\end{center}

\subsection{What ROV Can and Cannot Protect}

ROV is effective against many origin hijacks. If AS-161 announces
\texttt{10.154.0.0/24} while the ROA says only AS-154 is authorized, the route
is Invalid and can be rejected.

However, ROV does not protect the entire AS path. If the origin AS is valid but
the path contains false or manipulated intermediate ASes, origin validation
alone may still mark the route as Valid. ROV is therefore an important defense,
but it is not a complete solution to all BGP security problems.


% *******************************************
% SECTION
% *******************************************
\section{Lab Tasks}

% -------------------------------------------
% SUBSECTION
% -------------------------------------------
\subsection{Task 1: Getting Familiar with Route Origins}

The objective of this task is to identify the origin AS for routes inside the
emulator. Pick a BGP router outside AS-154, such as a transit AS connected to
AS-154's part of the topology. Get a shell on the router and inspect the route
to AS-154's prefix:

\begin{lstlisting}
# birdc show route all for 10.154.0.0/24
\end{lstlisting}

Please identify the following:

\begin{itemize}[noitemsep]
\item Which AS originates \texttt{10.154.0.0/24}?
\item What is the AS path?
\item Which BGP neighbor provided the route?
\item Can the router tell whether the origin AS is authorized?
\end{itemize}

In your report, explain why ordinary BGP route selection does not by itself
prove that the origin AS is legitimate.

\paragraph{Connection to theory.}
Relate your observation to the definition of the origin AS. Identify the last
AS in the AS path, and explain why this is the value checked by Route Origin
Validation.


% -------------------------------------------
% SUBSECTION
% -------------------------------------------
\subsection{Task 2: Launching a Prefix Hijacking Attack}

In this task, students will reproduce a basic origin hijack. AS-154 owns the
prefix \texttt{10.154.0.0/24}. We will make AS-161 announce the same prefix.
This simulates a malicious or misconfigured AS claiming reachability for a
prefix that it does not own.

On AS-161's BGP router, add the following static route to
\path{/etc/bird/bird.conf}. The exact location of this block depends on the
generated BIRD configuration, but it should be placed with the other static
routes.

\begin{lstlisting}
protocol static {
    ipv4 {
        table t_bgp;
    };
    route 10.154.0.0/24 blackhole {
        bgp_large_community.add(LOCAL_COMM);
    };
}
\end{lstlisting}

Reload BIRD:

\begin{lstlisting}
# birdc configure
\end{lstlisting}

From a different AS, such as AS-155 or AS-156, observe the route to
\texttt{10.154.0.0/24}. Use at least two of the following:

\begin{lstlisting}
# birdc show route all for 10.154.0.0/24
$ traceroute 10.154.0.71
$ ping 10.154.0.71
\end{lstlisting}

Please provide screenshots showing that AS-161's announcement changes the
selected route or disrupts reachability to AS-154. Explain whether this is an
origin hijack, a path manipulation attack, or both.

\paragraph{Connection to theory.}
Using the ROA model from the theory section, explain why AS-161's route for
\texttt{10.154.0.0/24} should be classified as Invalid.


% -------------------------------------------
% SUBSECTION
% -------------------------------------------
\subsection{Task 3: Creating a Small ROA Table}

The objective of this task is to understand what information a ROA contains. In
real RPKI, ROAs are signed objects published by prefix holders. In this lab, we
use a small text file to model the ROA database.

The provided \texttt{Labsetup/roa-table.txt} file contains entries like this:

\begin{lstlisting}
# prefix,maxLength,originAS,description
10.154.0.0/24,24,154,AS-154 customer prefix
10.155.0.0/24,24,155,AS-155 customer prefix
10.156.0.0/24,24,156,AS-156 customer prefix
10.161.0.0/24,24,161,AS-161 customer prefix
\end{lstlisting}

For each route below, decide whether the route is Valid, Invalid, or NotFound:

\begin{itemize}[noitemsep]
\item AS-154 announces \texttt{10.154.0.0/24}.
\item AS-161 announces \texttt{10.154.0.0/24}.
\item AS-154 announces \texttt{10.154.0.0/25}.
\item AS-170 announces \texttt{10.170.0.0/24}.
\end{itemize}

Please explain the role of \texttt{maxLength}. Why is
\texttt{10.154.0.0/25} not valid when the ROA says maxLength \texttt{24}?

\paragraph{Connection to theory.}
For each route, explicitly state which rule produced the validation result:
wrong origin AS, prefix length longer than maxLength, no covering ROA, or a
matching ROA.


% -------------------------------------------
% SUBSECTION
% -------------------------------------------
\subsection{Task 4: Installing an Educational ROV Filter}

In this task, students will configure a BGP router to reject invalid origins.
We will model ROV using a BIRD import filter. In a real deployment, the router
would normally receive validated prefix-origin data from an RPKI validator over
the RTR protocol. In this lab, we use a generated filter so students can focus
on the decision logic.

Generate the filter:

\begin{lstlisting}
$ cd Labsetup
$ python3 generate_bird_rov_filter.py roa-table.txt > seed_rov_filter.conf
\end{lstlisting}

Open \texttt{seed\_rov\_filter.conf}. It defines a function named
\texttt{seed\_rov\_invalid()}. Copy this function into the BIRD configuration
of AS-3's BGP router, or another provider router selected by the instructor.
Then call it from the import filter for the BGP session that receives routes
from AS-161:

\begin{lstlisting}
import filter {
    if seed_rov_invalid() then {
        print "ROV rejected invalid route: ", net, " origin ", bgp_path.last;
        reject;
    }
    accept;
};
\end{lstlisting}

Reload the BIRD configuration:

\begin{lstlisting}
# birdc configure
\end{lstlisting}

Now repeat the prefix hijacking attack from Task 2. Please answer the
following:

\begin{itemize}[noitemsep]
\item Does AS-3 still accept the fake route from AS-161?
\item Does the fake route propagate to other ASes?
\item What message or route-table evidence shows that the route was rejected?
\item Does AS-161 still have normal Internet connectivity after the filter is
      installed?
\end{itemize}

In your report, explain why this filter protects AS-154's prefix without
disconnecting AS-161 from the emulator.

\paragraph{Connection to theory.}
Map the generated BIRD filter to the theoretical ROV pipeline. Which part of
the filter models the ROA lookup? Which part models the Invalid decision? Which
part models the operator's policy to reject Invalid routes?


% -------------------------------------------
% SUBSECTION
% -------------------------------------------
\subsection{Task 5: Testing a More-Specific Prefix Attack}

Attackers often announce a more-specific prefix because routers prefer
longest-prefix matches during forwarding. In this task, AS-161 will announce
\texttt{10.154.0.0/25} instead of \texttt{10.154.0.0/24}.

On AS-161's BGP router, modify the malicious static route:

\begin{lstlisting}
protocol static {
    ipv4 {
        table t_bgp;
    };
    route 10.154.0.0/25 blackhole {
        bgp_large_community.add(LOCAL_COMM);
    };
}
\end{lstlisting}

Run the experiment twice:

\begin{itemize}[noitemsep]
\item First, run it without the ROV filter.
\item Second, run it with the ROV filter installed.
\end{itemize}

Please describe the difference. Explain why the ROA entry
\texttt{10.154.0.0/24,24,154} makes the \texttt{/25} announcement invalid,
even if AS-154 had originated it.

\paragraph{Connection to theory.}
Explain how longest-prefix forwarding makes more-specific hijacks powerful, and
how the ROA maxLength field limits which more-specific announcements can be
treated as Valid.


% -------------------------------------------
% SUBSECTION
% -------------------------------------------
\subsection{Task 6: Understanding the Limitations of ROV}

ROV checks whether the origin AS is authorized to originate a prefix. It does
not verify the entire AS path. This distinction is important.

Please answer the following questions:

\begin{itemize}[noitemsep]
\item What kind of hijack can ROV stop?
\item What kind of route manipulation can ROV fail to stop?
\item Why does a Valid origin not prove that the whole AS path is true?
\item Why might some operators choose to reject Invalid routes but still accept
      NotFound routes?
\end{itemize}

As an optional experiment, create a route announcement whose origin is valid but
whose path is less preferred or suspicious due to policy. Observe whether the
ROV filter rejects it. Explain the result.

\paragraph{Connection to theory.}
Use this task to separate origin validation from path validation. State exactly
which part of the BGP route ROV checks and which parts it leaves unchecked.


% *******************************************
% SECTION
% *******************************************
\section{Guidelines}

\paragraph{Finding the correct BGP session.}
Use \texttt{birdc show protocols} on the provider router to identify the BGP
session connected to AS-161. The protocol name often indicates the neighbor.
The emulator GUI can also help identify which ASes are connected.

\paragraph{Editing BIRD configuration.}
One way to edit the configuration is to copy it from the container to the host,
edit it, and copy it back:

\begin{lstlisting}
$ docker cp <container-id>:/etc/bird/bird.conf ./bird.conf

... edit bird.conf ...

$ docker cp ./bird.conf <container-id>:/etc/bird/bird.conf
$ docker exec <container-id> birdc configure
\end{lstlisting}

\paragraph{Valid, Invalid, and NotFound.}
Use the following rule in this lab:

\begin{itemize}[noitemsep]
\item Valid: a covering ROA exists, the origin AS matches, and the prefix
      length is no longer than maxLength.
\item Invalid: a covering ROA exists, but the origin AS or maxLength check
      fails.
\item NotFound: no covering ROA exists.
\end{itemize}

\paragraph{Instructor customization.}
Instructors can change the victim AS, attacker AS, provider AS, ROA table, and
maxLength values. A useful variation is to give students different victim
prefixes or to intentionally omit one ROA and ask students to explain why the
route becomes NotFound instead of Invalid.


% *******************************************
% SECTION
% *******************************************
\section{Submission}

%%%%%%%%%%%%%%%%%%%%%%%%%%%%%%%%%%%%%%%%
\input{\commonfolder/submission}
%%%%%%%%%%%%%%%%%%%%%%%%%%%%%%%%%%%%%%%%

In addition, your report should include:

\begin{itemize}[noitemsep]
\item Screenshots of route tables before and after the hijack.
\item The ROA table used in your experiment.
\item The BIRD ROV filter installed on the provider router.
\item Evidence that invalid routes are rejected.
\item An explanation of the limitations of Route Origin Validation.
\end{itemize}

\end{document}
